
\documentclass[12pt,a4paper]{article}

\usepackage[T2A]{fontenc}
\usepackage[utf8]{inputenc}
\usepackage[russian]{babel}
\usepackage{amsmath,amssymb,amsfonts}
\usepackage{geometry}
\usepackage{graphicx}
\usepackage{booktabs}
\usepackage{array}
\usepackage{float}
\usepackage{caption}
\usepackage{subcaption}
\usepackage{siunitx}
\usepackage{enumitem}
\usepackage{hyperref}
\usepackage{xcolor}

\geometry{
    left=25mm,
    right=20mm,
    top=20mm,
    bottom=20mm
}

\hypersetup{
    colorlinks=true,
    linkcolor=black,
    urlcolor=blue,
    citecolor=black
}

\newcommand{\vect}[1]{\boldsymbol{#1}}
\newcommand{\matr}[1]{\mathbf{#1}}

% ------------------------------------------------------------
% Универсальное место под рисунок.
% Вместо \figureplaceholder можно позднее поставить \includegraphics.
% ------------------------------------------------------------
\newcommand{\figureplaceholder}[2][0.88\textwidth]{%
\begin{figure}[H]
    \centering
    \fbox{%
        \begin{minipage}[c][0.26\textheight][c]{#1}
            \centering
            \textbf{Место для рисунка}\\[1ex]
            #2
        \end{minipage}
    }
\end{figure}
}

\title{
    Численное моделирование электромагнитных задач\\
    методом частичных эквивалентных схем
}
\author{}
\date{}

\begin{document}

\maketitle
\tableofcontents
\newpage


% ============================================================
\section{Введение}
% ============================================================

В работе реализован численный решатель на основе метода частичных
эквивалентных схем (ЧЭС, PEEC --- Partial Element Equivalent Circuit).
Метод применяется для решения как частотных, так и нестационарных
электромагнитных задач на поверхностных сетках, состоящих из
четырёхугольных элементов.

В качестве основных тестовых и прикладных задач рассмотрены:

\begin{enumerate}[label=\arabic*)]
    \item рассеивание падающей электромагнитной волны на квадратной пластине;
    \item расчёт эффективной площади рассеяния квадратной пластины;
    \item нестационарная задача удара молнии в проводящий корпус с апертурой
    и внутренней электрической цепью.
\end{enumerate}

Основными неизвестными являются токи, связанные с рёбрами дискретизации,
и узловые электрические величины. Для нестационарной задачи используется
неявная аппроксимация первого или второго порядка по времени.


% ============================================================
\section{Численная схема метода ЧЭС}
% ============================================================

\subsection{Дискретизация поверхности}

Проводящая поверхность разбивается на четырёхугольные элементы.
Глобальные токовые неизвестные $I_i$ связываются с рёбрами сетки.
Для ребра $i$ направление определяется в момент его первого появления
при обходе элементов.

Для каждого ребра вводится ориентированный единичный вектор
$\vect e_i$. Плотность поверхностного тока в соответствующей области
можно представить в виде
\begin{equation}
    \vect J_i
    =
    \frac{I_i}{w_i}\vect e_i ,
\end{equation}
где $w_i$ --- эффективная ширина области, относящейся к ребру.

Узловые электрические величины определяются на двойственной сетке,
построенной вокруг узлов исходной поверхностной сетки.

\figureplaceholder{
Схема поверхностной сетки, ребровых токов и двойственных областей узлов.
}


\subsection{Матрица инцидентности}

Связь между рёбрами и узлами задаётся матрицей инцидентности
\begin{equation}
    \matr A \in \mathbb R^{N_e\times N_v}.
\end{equation}

Каждому глобальному ребру назначается ориентация
\begin{equation}
    e=(a\rightarrow b).
\end{equation}
В используемой реализации эта ориентация определяется при первом появлении
ребра в процессе обхода четырёхугольных элементов сетки.

Строка матрицы $\matr A$, соответствующая ребру $e$, содержит только
два ненулевых элемента:
\begin{equation}
    A_{e,a}=-1,
    \qquad
    A_{e,b}=+1.
\end{equation}

Все остальные элементы этой строки равны нулю:
\begin{equation}
    A_{e,j}=0,
    \qquad
    j\neq a,b.
\end{equation}

Таким образом, матрица $\matr A$ является чисто топологическим оператором
и не зависит от геометрических размеров элементов. Оператор
$\matr A^T$ переводит ребровые токи в узловые токовые балансы:
\begin{equation}
    \left(\matr A^T\vect I\right)_j
    =
    \sum_{e=1}^{N_e}
    A_{e,j}I_e.
\end{equation}

Именно этот оператор входит в дискретный закон сохранения заряда.


\subsection{Матрица частичных индуктивностей}

Электромагнитная связь между токовыми областями учитывается
матрицей частичных индуктивностей
\begin{equation}
    \matr L\in\mathbb R^{N_e\times N_e}.
\end{equation}

Для пары токовых областей $a$ и $b$ используется интегральное выражение
\begin{equation}
    L_{ab}
    =
    \frac{\mu_0}{4\pi}
    \left(
        \vect e_a\cdot\vect e_b
    \right)
    \iint_{\Pi_a^e\times\Pi_b^e}
    \frac{
        \mathrm d r\,\mathrm d r'
    }{
        |\vect r-\vect r'|
    },
    \label{eq:L_integral}
\end{equation}
где $\Pi_a^e$ и $\Pi_b^e$ --- области поверхности, относящиеся
к соответствующим рёбрам, а $\vect e_a$, $\vect e_b$ --- единичные
векторы направлений токов.

В дискретной модели каждая ребровая область строится как сумма частей
соседних четырёхугольных ячеек:
\begin{equation}
    \Pi_i^e
    =
    P_{1i}+P_{2i}.
\end{equation}

Диагональные элементы
\begin{equation}
    L_{aa}
\end{equation}
задают собственную частичную индуктивность ребровой области,
а внедиагональные элементы
\begin{equation}
    L_{ab},\qquad a\neq b,
\end{equation}
описывают взаимную магнитную связь.

Интеграл~\eqref{eq:L_integral} вычисляется численно квадратурой
по соответствующим поверхностным областям. Для регулярных пар областей
используется обычное численное интегрирование. Для собственных и близких
взаимодействий необходимо повышать точность квадратуры, поскольку ядро
\begin{equation}
    \frac{1}{|\vect r-\vect r'|}
\end{equation}
становится быстро меняющимся или слабо сингулярным.

После вычисления всех парных взаимодействий формируется плотная матрица
\begin{equation}
    \matr L=
    \begin{bmatrix}
        L_{11} & \cdots & L_{1N_e}\\
        \vdots & \ddots & \vdots\\
        L_{N_e1} & \cdots & L_{N_eN_e}
    \end{bmatrix}.
\end{equation}


\subsection{Матрица коэффициентов потенциала}

Электростатическая связь между двойственными узловыми областями
описывается матрицей коэффициентов потенциала
\begin{equation}
    \matr P\in\mathbb R^{N_v\times N_v}.
\end{equation}

Каждому узлу $j$ ставится в соответствие двойственная область
$\Pi_j^v$. Для четырёхугольной сетки она собирается из локальных
частей соседних ячеек:
\begin{equation}
    \Pi_j^v
    =
    \sum_k q_{jk},
\end{equation}
где $q_{jk}$ --- часть $k$-й ячейки, относящаяся к узлу $j$.

Коэффициент потенциала между двумя узловыми областями можно записать
в виде
\begin{equation}
    P_{ij}
    =
    \frac{1}{4\pi\varepsilon_0}
    \frac{1}{
        |\Pi_i^v|\,
        |\Pi_j^v|
    }
    \iint_{\Pi_i^v\times\Pi_j^v}
    \frac{
        \mathrm dS\,\mathrm dS'
    }{
        |\vect r-\vect r'|
    }.
    \label{eq:P_integral}
\end{equation}

Формула~\eqref{eq:P_integral} соответствует модели с постоянной
плотностью заряда внутри каждой двойственной области. В общем случае
масштабирующий множитель может быть включён в определение базисной
функции, однако в программе сохраняется физический смысл
\begin{equation}
    \vect\varphi
    =
    \matr P\vect q.
\end{equation}

Матрица $\matr P$ имеет размерность
\begin{equation}
    [P_{ij}]=\si{\per\farad}.
\end{equation}

Как и в случае $\matr L$, интегралы вычисляются численно квадратурой
по двойственным поверхностным областям. Особое внимание уделяется
диагональным элементам $P_{ii}$, поскольку они содержат собственное
электростатическое взаимодействие области.

Для перехода к используемой в нестационарном решателе форме вводится
диагональная матрица
\begin{equation}
    \matr F
    =
    \operatorname{diag}
    \left(
        F_1,\ldots,F_{N_v}
    \right),
\end{equation}
где
\begin{equation}
    F_j
    =
    \frac{1}{P_{jj}}.
\end{equation}

Внутренняя узловая неизвестная определяется как собственное напряжение
\begin{equation}
    V_{c,j}
    =
    P_{jj}q_j.
\end{equation}
Отсюда непосредственно
\begin{equation}
    q_j
    =
    F_jV_{c,j}.
\end{equation}

Для восстановления полного потенциала строится матрица
\begin{equation}
    \matr S,
\end{equation}
элементы которой равны
\begin{equation}
    S_{ja}
    =
    \frac{P_{ja}}{P_{aa}}.
    \label{eq:S_definition}
\end{equation}

Действительно,
\begin{align}
    \varphi_j
    &=
    \sum_{a=1}^{N_v}
    P_{ja}q_a
    =
    \sum_{a=1}^{N_v}
    P_{ja}
    \frac{V_{c,a}}{P_{aa}}
    \nonumber\\
    &=
    \sum_{a=1}^{N_v}
    S_{ja}V_{c,a}.
\end{align}

Следовательно,
\begin{equation}
    \vect\varphi
    =
    \matr S\vect V_c,
\end{equation}
и одновременно
\begin{equation}
    \vect q
    =
    \matr F\vect V_c.
\end{equation}

Для ребрового уравнения заранее формируется оператор
\begin{equation}
    \matr{AS}
    =
    \matr A\matr S.
\end{equation}

В программе его можно вычислять непосредственно, не формируя
промежуточные операции в каждом временном шаге:
\begin{equation}
    (AS)_{e,a}
    =
    \frac{
        P_{n,a}-P_{m,a}
    }{
        P_{aa}
    },
\end{equation}
где $m$ и $n$ --- начальный и конечный узлы ориентированного ребра $e$.


\subsection{Сопротивления}

Омические потери учитываются диагональной матрицей
\begin{equation}
    \matr R.
\end{equation}

Для идеально проводящего тела (PEC)
\begin{equation}
    \matr R = \matr 0.
\end{equation}


% ============================================================
\subsection{Непрерывная система уравнений}
% ============================================================

В используемой форме нестационарная система ЧЭС записывается как
\begin{equation}
    \matr L
    \frac{\mathrm d\vect I}{\mathrm dt}
    +
    \matr R\vect I
    +
    \matr A\matr S\vect V_c
    +
    \vect U_e
    =
    \vect 0,
    \label{eq:edge_continuous}
\end{equation}
\begin{equation}
    -\matr A^T\vect I
    +
    \matr F
    \frac{\mathrm d\vect V_c}{\mathrm dt}
    =
    \vect J,
    \label{eq:node_continuous}
\end{equation}
где $\vect U_e$ --- внешнее возбуждение в ребровых уравнениях,
а $\vect J$ --- узловой токовый источник.


% ============================================================
\subsection{Неявная аппроксимация по времени}
% ============================================================

Для временной аппроксимации используется обобщённая $\theta$-схема.

При
\begin{equation}
    \theta=1
\end{equation}
получается схема backward Euler первого порядка.

При
\begin{equation}
    \theta=\frac12
\end{equation}
получается неявная трапецеидальная схема второго порядка.

Для шага $\Delta t$ левая часть линейной системы имеет вид
\begin{equation}
\begin{bmatrix}
\dfrac{\matr L}{\Delta t}
+\theta\matr R
&
\theta\matr A\matr S
\\[2ex]
-\theta\matr A^T
&
\dfrac{\matr F}{\Delta t}
\end{bmatrix}
\begin{bmatrix}
\vect I^{n+1}
\\
\vect V_c^{n+1}
\end{bmatrix}
=
\begin{bmatrix}
\vect b_I
\\
\vect b_V
\end{bmatrix}.
\label{eq:theta_system}
\end{equation}

Правая часть ребрового уравнения:
\begin{align}
\vect b_I
={}&
\frac{\matr L}{\Delta t}\vect I^n
-
(1-\theta)\matr R\vect I^n
-
(1-\theta)\matr A\matr S\vect V_c^n
\nonumber\\
&-
(1-\theta)\vect U_e^n
-
\theta\vect U_e^{n+1}.
\end{align}

Правая часть узлового уравнения:
\begin{align}
\vect b_V
={}&
\frac{\matr F}{\Delta t}\vect V_c^n
+
(1-\theta)\matr A^T\vect I^n
\nonumber\\
&+
(1-\theta)\vect J^n
+
\theta\vect J^{n+1}.
\end{align}

Матрица системы при постоянном $\Delta t$ остаётся неизменной,
поэтому её LU-разложение выполняется один раз до начала временного цикла.

Практический алгоритм одного шага имеет вид:
\begin{enumerate}[label=\arabic*.]
    \item По известным $\vect I^n$ и $\vect V_c^n$ формируются
    значения внешних воздействий $\vect J^n$, $\vect J^{n+1}$,
    $\vect U_e^n$, $\vect U_e^{n+1}$.

    \item По формулам $\theta$-схемы собирается правая часть
    \begin{equation}
        \vect b^{n+1}.
    \end{equation}

    \item Решается линейная система
    \begin{equation}
        \matr K_\theta
        \vect x^{n+1}
        =
        \vect b^{n+1},
    \end{equation}
    где
    \begin{equation}
        \vect x^{n+1}
        =
        \begin{bmatrix}
            \vect I^{n+1}\\
            \vect V_c^{n+1}
        \end{bmatrix}.
    \end{equation}

    \item Из $\vect V_c^{n+1}$ восстанавливаются физические величины:
    \begin{equation}
        \vect q^{n+1}
        =
        \matr F\vect V_c^{n+1},
    \end{equation}
    \begin{equation}
        \vect\varphi^{n+1}
        =
        \matr S\vect V_c^{n+1}.
    \end{equation}

    \item Выполняется переход
    \begin{equation}
        n\leftarrow n+1.
    \end{equation}
\end{enumerate}

Для постоянного шага времени матрица
\begin{equation}
    \matr K_\theta
\end{equation}
не изменяется, поэтому вычислительно наиболее дорогая операция
LU-факторизации производится только один раз. На каждом следующем шаге
выполняются только сборка правой части и обратная подстановка.

Выбор $\Delta t$ определяется характерным временем изменения источника
и собственными временными масштабами PEEC-системы. Для импульса молнии
необходимо, чтобы быстрый фронт двойной экспоненты был разрешён несколькими
временными шагами. В выполненных расчётах использовалось
\begin{equation}
    \Delta t=10^{-7}\ \si{\second}.
\end{equation}


% ============================================================
\subsection{Частотная постановка}
% ============================================================

При гармонической зависимости
\begin{equation}
    e^{j\omega t}
\end{equation}
производные заменяются умножением на $j\omega$.

Для задачи рассеивания получается система
\begin{equation}
\begin{bmatrix}
\matr R + j\omega\matr L
&
\matr A\matr S
\\
-\matr A^T
&
j\omega\matr F
\end{bmatrix}
\begin{bmatrix}
\vect I
\\
\vect V_c
\end{bmatrix}
=
\begin{bmatrix}
-\vect U_{\mathrm{inc}}
\\
\vect 0
\end{bmatrix}.
\label{eq:harmonic_system}
\end{equation}

Вектор $\vect U_{\mathrm{inc}}$ формируется интегрированием
падающего электрического поля вдоль рёбер сетки.


% ============================================================
\subsection{Внутренняя резистивная ветвь}
% ============================================================

Для внутренней резистивной нагрузки между узлами $A$ и $B$
вводится вектор
\begin{equation}
    b_R[A]=-1,
    \qquad
    b_R[B]=+1.
\end{equation}

Физическое напряжение ветви:
\begin{equation}
    b_R^T\vect\varphi
    =
    b_R^T\matr S\vect V_c.
\end{equation}

Уравнение резистора:
\begin{equation}
    b_R^T\matr S\vect V_c
    +
    R_{\mathrm{sh}}I_{\mathrm{sh}}
    =
    0.
\end{equation}


% ============================================================
\subsection{Эквивалентные ячейки апертуры}
% ============================================================

Для сокращённой одноэтапной модели апертура представляется
набором распределённых эквивалентных ячеек.
Каждая ячейка $m$ соединяет противоположные узлы апертуры
$a_m$ и $b_m$.

Для неё вводится
\begin{equation}
    b_m[a_m]=-1,
    \qquad
    b_m[b_m]=+1.
\end{equation}

Параметры локальной ячейки:
\begin{equation}
    L_m=L'\Delta l_m,
    \qquad
    C_m=C'\Delta l_m,
\end{equation}
где $\Delta l_m$ --- продольная длина поддержки ячейки.

Для линии с характеристическим сопротивлением $Z_{\mathrm{slot}}$
и скоростью распространения $v_{\mathrm{slot}}$
\begin{equation}
    L'
    =
    \frac{Z_{\mathrm{slot}}}{v_{\mathrm{slot}}},
    \qquad
    C'
    =
    \frac{1}{
        Z_{\mathrm{slot}}v_{\mathrm{slot}}
    }.
\end{equation}

Ёмкостные ветви модифицируют узловой оператор:
\begin{equation}
    \matr F_{\mathrm{slot}}
    =
    \matr F
    +
    \sum_{m=1}^{N_s}
    C_m
    b_m
    c_m^T,
\end{equation}
где
\begin{equation}
    c_m^T
    =
    b_m^T\matr S.
\end{equation}

В результате число дополнительных неизвестных равно числу
индуктивных токов эквивалентных ячеек $N_s$, что существенно меньше,
чем при решении второй полной PEEC-системы.

\figureplaceholder{
Схема апертуры, заменённой распределёнными эквивалентными LC-ячейками.
}


% ============================================================
\section{Результаты расчётов}
% ============================================================


% ============================================================
\subsection{Рассеивание падающей волны на квадратной пластине}
% ============================================================

В первой тестовой задаче рассматривается квадратная идеально проводящая
пластина, на которую падает гармоническая плоская электромагнитная волна.

В расчёте задаются:
\begin{equation}
    f = \text{[указать частоту]},
\end{equation}
\begin{equation}
    E_0 = \text{[указать амплитуду]},
\end{equation}
а также углы падения $\theta$, $\varphi$ и поляризация волны.

После построения матриц $\matr L$, $\matr P$, $\matr R$,
$\matr A$ и формирования вектора падающего поля решается
гармоническая система~\eqref{eq:harmonic_system}.

Основным результатом является комплексное распределение поверхностных
токов по пластине.

\figureplaceholder{
Геометрия квадратной пластины и направление падающей электромагнитной волны.
}

\figureplaceholder{
Распределение модуля поверхностного тока $|\vect J|$ на квадратной пластине.
}

\figureplaceholder{
Распределение фазы или действительной части поверхностного тока на пластине.
}

\begin{table}[H]
    \centering
    \caption{Основные параметры задачи рассеивания}
    \begin{tabular}{>{\raggedright\arraybackslash}p{0.48\textwidth}
                    >{\centering\arraybackslash}p{0.34\textwidth}}
        \toprule
        Параметр & Значение \\
        \midrule
        Размер стороны пластины & [вставить] \\
        Частота $f$ & [вставить] \\
        Амплитуда $E_0$ & [вставить] \\
        Угол $\theta$ & [вставить] \\
        Угол $\varphi$ & [вставить] \\
        Поляризация & [вставить] \\
        Число узлов & [вставить] \\
        Число рёбер & [вставить] \\
        Число элементов & [вставить] \\
        \bottomrule
    \end{tabular}
\end{table}

Полученное распределение токов может использоваться далее
для восстановления рассеянного поля и расчёта ЭПР.


% ============================================================
\subsection{Эффективная площадь рассеяния квадратной пластины}
% ============================================================

Эффективная площадь рассеяния определяется через дальнее
рассеянное поле:
\begin{equation}
    \sigma
    =
    \lim_{r\rightarrow\infty}
    4\pi r^2
    \frac{
        |\vect E_{\mathrm{sca}}|^2
    }{
        |\vect E_{\mathrm{inc}}|^2
    }.
\end{equation}

После решения гармонической задачи вычисляется поле,
создаваемое найденными поверхностными токами в заданном
направлении наблюдения.

Расчёт выполняется для набора углов наблюдения,
например при фиксированном $\theta_{\mathrm{obs}}$
и изменении азимутального угла $\varphi$.

\figureplaceholder{
Схема определения углов наблюдения при расчёте ЭПР квадратной пластины.
}

\figureplaceholder{
ЭПР квадратной пластины $\sigma(\varphi)$ в линейном масштабе.
}

\figureplaceholder{
ЭПР квадратной пластины в логарифмическом масштабе, дБм$^2$.
}

\begin{table}[H]
    \centering
    \caption{Параметры расчёта ЭПР}
    \begin{tabular}{>{\raggedright\arraybackslash}p{0.48\textwidth}
                    >{\centering\arraybackslash}p{0.34\textwidth}}
        \toprule
        Параметр & Значение \\
        \midrule
        Электрический размер $ka$ & [вставить] \\
        Характерный размер $a$ & [вставить] \\
        Частота & [вставить] \\
        $\theta_{\mathrm{obs}}$ & [вставить] \\
        Диапазон $\varphi$ & [вставить] \\
        Шаг по $\varphi$ & [вставить] \\
        Максимальная ЭПР & [вставить] \\
        Угол максимума & [вставить] \\
        \bottomrule
    \end{tabular}
\end{table}


% ============================================================
\subsection{Нестационарная задача удара молнии}
% ============================================================

Для моделирования тока молнии используется двойная экспонента
\begin{equation}
    I_L(t)
    =
    K
    \left[
        e^{-\alpha(t-t_d)}
        -
        e^{-\beta(t-t_d)}
    \right],
    \qquad
    t\ge t_d.
\end{equation}

В использованных расчётах применялись параметры
\begin{equation}
    K = 102900~\si{\ampere},
\end{equation}
\begin{equation}
    \alpha = 1500~\si{\per\second},
    \qquad
    \beta = 10^6~\si{\per\second}.
\end{equation}

Для этих параметров максимум импульса достигается приблизительно
при
\begin{equation}
    t_{\mathrm{peak}}
    \approx
    6.51~\si{\micro\second}.
\end{equation}

\figureplaceholder{
График задающего тока молнии $I_L(t)$.
}


% ============================================================
\subsubsection{Последовательность решения двухэтапной задачи}
% ============================================================

Для корпуса с апертурой реализована двухэтапная последовательность решения.

\begin{enumerate}[label=\arabic*.]
    \item Формируется закрытая модель корпуса, в которой апертура временно
    перекрывается проводящей поверхностью.

    \item Для закрытой модели собираются PEEC-матрицы
    $\matr L$, $\matr P$, $\matr R$, $\matr A$.

    \item В выбранный узел прикладывается ток молнии, а во втором узле
    задаётся возвратный ток.

    \item На каждом временном шаге решается нестационарная PEEC-система
    закрытого корпуса.

    \item По токам рёбер временной крышки определяется токовый перенос
    через границу апертуры:
    \begin{equation}
        \vect I_\Gamma
        =
        \matr B_\Gamma
        \vect I_{\mathrm{cover}}.
    \end{equation}

    \item Формируется открытая модель корпуса с реальной апертурой.

    \item Вектор $\vect I_\Gamma$ используется как узловой источник
    второй PEEC-системы.

    \item При наличии внутреннего резистора одновременно вычисляется
    ток $I_{\mathrm{shunt}}(t)$ и напряжение между его выводами.

    \item Для каждого временного шага сохраняются поверхностные токи,
    узловые потенциалы и ток внутренней ветви.
\end{enumerate}

\figureplaceholder{
Блок-схема двухэтапного алгоритма:
закрытая модель $\rightarrow$ ток через апертуру $\rightarrow$ открытая модель.
}

\figureplaceholder{
Геометрия закрытой модели корпуса с временной PEC-крышкой.
}

\figureplaceholder{
Геометрия открытой модели корпуса с апертурой.
}

\figureplaceholder{
Распределение поверхностного тока на закрытом корпусе в момент,
близкий к максимуму импульса молнии.
}

\figureplaceholder{
Распределение поверхностного тока на открытом корпусе после переноса
возбуждения через апертуру.
}

\figureplaceholder{
График тока внутреннего шунта $|I_{\mathrm{shunt}}(t)|$.
}


% ============================================================
\subsubsection{Одноэтапная модель с эквивалентными ячейками щели}
% ============================================================

В качестве сокращённой альтернативы двухэтапному решению реализована
модель, в которой реальная апертура заменяется распределённым набором
эквивалентных LC-ячеек.

Для каждой пары противоположных узлов апертуры автоматически
определяются внутренние PEEC-индексы и локальная длина поддержки
$\Delta l_m$. После этого параметры ячейки вычисляются как
\begin{equation}
    L_m=L'\Delta l_m,
    \qquad
    C_m=C'\Delta l_m.
\end{equation}

При этом используется только одна открытая PEEC-сетка,
а неизвестные имеют вид
\begin{equation}
    \vect x
    =
    \begin{bmatrix}
        \vect I \\
        \vect V_c \\
        \vect I_{\mathrm{slot}} \\
        I_{\mathrm{shunt}}
    \end{bmatrix}.
\end{equation}

Такой подход уменьшает число дополнительных неизвестных
по сравнению со второй полной PEEC-системой.

\figureplaceholder{
Разбиение апертуры на распределённые эквивалентные LC-ячейки.
}

\figureplaceholder{
Графики токов отдельных эквивалентных ячеек
$I_{\mathrm{slot},m}(t)$.
}

\figureplaceholder{
Огибающая $\max_m |I_{\mathrm{slot},m}(t)|$.
}

\figureplaceholder{
Сравнение токов внутреннего шунта для двухэтапной модели
и модели с эквивалентными ячейками щели.
}


% ============================================================
\subsubsection{Параметры временного расчёта}
% ============================================================

\begin{table}[H]
    \centering
    \caption{Параметры расчёта удара молнии}
    \begin{tabular}{>{\raggedright\arraybackslash}p{0.48\textwidth}
                    >{\centering\arraybackslash}p{0.34\textwidth}}
        \toprule
        Параметр & Значение \\
        \midrule
        $K$ & $102900~\si{\ampere}$ \\
        $\alpha$ & $1500~\si{\per\second}$ \\
        $\beta$ & $10^6~\si{\per\second}$ \\
        Шаг времени $\Delta t$ & $10^{-7}~\si{\second}$ \\
        Конечное время & $10^{-4}~\si{\second}$ \\
        Порядок временной схемы & [1 или 2] \\
        Сопротивление внутреннего шунта & $1000~\si{\ohm}$ \\
        Число эквивалентных ячеек апертуры & [вставить] \\
        \bottomrule
    \end{tabular}
\end{table}


% ============================================================
\section{Обсуждение результатов}
% ============================================================

Проведённые расчёты позволяют проверить работу одной и той же
ЧЭС-дискретизации в частотной и временной областях.

Для квадратной пластины гармонический расчёт даёт распределение
индуцированных поверхностных токов, которое далее используется
при восстановлении рассеянного поля и вычислении ЭПР.

Для нестационарной задачи молнии двухэтапная постановка позволяет
явно разделить внешнюю задачу возбуждения закрытого корпуса
и проникновение тока через апертуру во внутреннюю область.

Одноэтапная модель с эквивалентными ячейками предназначена
для сокращения размера задачи. Её точность должна оцениваться
сравнением с двухэтапным расчётом по таким величинам, как
$I_{\mathrm{shunt}}(t)$, напряжение между внутренними выводами,
максимальные поверхностные токи и распределение токов вдоль апертуры.

\figureplaceholder{
Итоговый сравнительный рисунок основных результатов расчёта.
}


% ============================================================
\section{Заключение}
% ============================================================

Разработан численный решатель ЧЭС для решения гармонических
и нестационарных электромагнитных задач на поверхностных сетках.

Реализованы:

\begin{enumerate}[label=\arabic*)]
    \item формирование топологии и матрицы инцидентности;
    \item вычисление матриц частичных индуктивностей и коэффициентов потенциала;
    \item гармоническое решение задачи рассеивания;
    \item восстановление дальнего рассеянного поля и расчёт ЭПР;
    \item неявное интегрирование нестационарной PEEC-системы;
    \item двухэтапная задача удара молнии в корпус с апертурой;
    \item внутренняя резистивная ветвь;
    \item одноэтапная модель апертуры с распределёнными
    эквивалентными LC-ячейками;
    \item экспорт результатов для последующей визуализации.
\end{enumerate}

Полученная программная структура позволяет далее проводить
сравнение полной двухэтапной модели с сокращённой моделью апертуры
и исследовать влияние параметров сетки, шага времени и числа
эквивалентных ячеек на точность решения.


\end{document}
